% arara: pdflatex: {shell: yes, files: [latexindent]}
\section{Fine tuning}\label{sec:finetuning}
\fancyhead[R]{\bfseries\thepage%
	\tikz[remember picture,overlay] {
		\node at (1,0){\includegraphics{logo-bw}}; }
}
\texttt{latexindent.pl} operates by looking for the code blocks detailed in \vref{tab:code-blocks}.
\announce{new}{details of fine tuning of code blocks} The fine tuning of the details of such code blocks
is controlled by the \texttt{fineTuning} field, detailed in \cref{lst:fineTuning}.

This field is for those that would like to peek under the bonnet/hood and make some
fine tuning to \texttt{latexindent.pl}'s operating.

\begin{warning}
	Making changes to the fine tuning may have significant consequences for your indentation scheme,
	proceed with caution!
\end{warning}

\begin{widepage}
\cmhlistingsfromfile*[style=fineTuning]*{../defaultSettings.yaml}[width=0.95\linewidth,before=\centering,yaml-TCB]{\texttt{fineTuning}}{lst:fineTuning}
\end{widepage}

The fields given in \cref{lst:fineTuning} are all \emph{regular expressions}. This manual is not intended to be
a tutorial on regular expressions; you might like to read, for example, \cite{masteringregexp} for a detailed
covering of the topic.

We make the following comments with reference to \cref{lst:fineTuning}:
\begin{enumerate}
	\item the \texttt{environments:name} field details that the \emph{name} of an environment can contain:
	      \begin{enumerate}
		      \item \texttt{a-z} lower case letters
		      \item \texttt{A-Z} upper case letters
		      \item \texttt{@} the \texttt{@} 'letter'
		      \item \lstinline!\*! stars
		      \item \texttt{0-9} numbers
		      \item \lstinline!_! underscores
		      \item \lstinline!\! backslashes
	      \end{enumerate}
	      The \texttt{+} at the end means \emph{at least one} of the above characters.
	\item the \texttt{ifElseFi:name} field:
	      \begin{enumerate}
		      \item \lstinline^@?^ means that it \emph{can possibly} begin with \lstinline^@^
		      \item followed by \texttt{if}
		      \item followed by 0 or more characters from \texttt{a-z}, \texttt{A-Z} and \texttt{@}
              \item the \texttt{?} the end means \emph{non-greedy}, which means `stop the match as soon as possible'
	      \end{enumerate}
	\item the \texttt{keyEqualsValuesBracesBrackets} contains some interesting syntax:
	      \begin{enumerate}
		      \item \lstinline!|! means `or'
		      \item \lstinline^(?:(?<!\\)\{)^ the \lstinline^(?:...)^ uses a \emph{non-capturing} group -- you don't necessarily need to worry about
		            what this means, but just know that for the \texttt{fineTuning} feature you should only ever use
		            \emph{non}-capturing groups, and \emph{not} capturing groups, which are simply \lstinline!(...)!
		      \item \lstinline^(?<!\\)\{)^ means a \lstinline^{^ but it can \emph{not} be immediately preceeded by a \lstinline!\!
	      \end{enumerate}
	\item in the \texttt{arguments:before} field
	      \begin{enumerate}
            \item \lstinline^\d\h*^ means a digit (i.e. a number), followed by 0 or more horizontal spaces
		      \item \lstinline^;?,?^ means \emph{possibly} a semi-colon, and possibly a comma
		      \item \lstinline^\<.*?\>^ is designed for 'beamer'-type commands; the \lstinline^.*?^ means anything in between \lstinline^<...>^
	      \end{enumerate}
        \item the \texttt{modifyLineBreaks} field refers to fine tuning settings detailed in \vref{sec:modifylinebreaks}. In particular:
	      \begin{enumerate}
            \item \texttt{betterFullStop} is in relation to the one sentence per line routine, detailed in \vref{sec:onesentenceperline}
            \item \texttt{doubleBackSlash} is in relation to the \texttt{DBSStartsOnOwnLine} and \texttt{DBSFinishesWithLineBreak} polyswitches surrounding double back slashes, see \vref{subsec:dbs}
            \item \texttt{comma} is in relation to the \texttt{CommaStartsOnOwnLine} and \texttt{CommaFinishesWithLineBreak} polyswitches surrounding commas in optional and mandatory arguments; see \vref{tab:poly-switch-mapping}
	      \end{enumerate}
\end{enumerate}

It is not obvious from \cref{lst:fineTuning}, but each of the \texttt{follow}, \texttt{before} and \texttt{between} fields
allow trailing comments, line breaks, and horizontal spaces between each character.

As a demonstration, consider the file given in \cref{lst:finetuning1}, together with its default output using the command
\begin{commandshell}
latexindent.pl finetuning1.tex 
\end{commandshell}
is given in \cref{lst:finetuning1-default}.

        \begin{cmhtcbraster}[raster column skip=.01\linewidth]
		\cmhlistingsfromfile*{demonstrations/finetuning1.tex}{\texttt{finetuning1.tex}}{lst:finetuning1}
		\cmhlistingsfromfile*{demonstrations/finetuning1-default.tex}{\texttt{finetuning1.tex} default}{lst:finetuning1-default}
        \end{cmhtcbraster}

It's clear from \cref{lst:finetuning1-default} that the indentation scheme has not worked as expected. We can \emph{fine tune} the 
indentation scheme by employing the settings given in \cref{lst:fine-tuning1} and running the command
\begin{commandshell}
latexindent.pl finetuning1.tex -l=fine-tuning1.yaml
\end{commandshell}
and the associated (desired) output is given in \cref{lst:finetuning1-mod1}.

        \begin{cmhtcbraster}[raster column skip=.01\linewidth]
		\cmhlistingsfromfile*{demonstrations/finetuning1-mod1.tex}{\texttt{finetuning1.tex} using \cref{lst:fine-tuning1}}{lst:finetuning1-mod1}
		\cmhlistingsfromfile*[style=yaml-LST]*{demonstrations/fine-tuning1.yaml}[yaml-TCB]{\texttt{finetuning1.yaml}}{lst:fine-tuning1}
        \end{cmhtcbraster}
